\documentclass[11pt]{article}
\usepackage[margin=1in]{geometry}
\usepackage{amsmath,amssymb,amsthm,mathtools}
\usepackage[hidelinks,hypertexnames=false]{hyperref}
\usepackage{xcolor}
\usepackage{enumitem}
\newtheorem{theorem}{Theorem}
\newcommand{\Hcal}{\mathcal H}
\newcommand{\Tr}{\operatorname{Tr}}
\newcommand{\MMD}{\operatorname{MMD}}
\setlist{nosep}
\title{Minimax-rate optimality of MONK under adversarial contamination\\
\large A literature-implied answer to the question in arXiv:1802.04784}
\author{Autonomous mathematical-discovery run \texttt{fa\_banach\_001}}
\date{12 August 2026}

\begin{document}
\maketitle

\noindent\fcolorbox{black}{gray!8}{\parbox{0.94\linewidth}{
\textbf{Status: literature-implied answer (full minimax-rate resolution).}
The order of MONK's error due to adversarial outliers is optimal under the
paper's finite-second-moment hypothesis.  The relation to MONK is identified
here; the supporting authors do not explicitly state an RKHS corollary.  Exact
leading constants and the exact breakdown behavior as the fraction of
contaminated blocks approaches one half are not claimed.
}}

\section*{Source question and supporting result}

The source is M. Lerasle, Z. Szab\'o, T. Mathieu, and G. Lecu\'e,
\emph{MONK---Outlier-Robust Mean Embedding Estimation by
Median-of-Means}, arXiv:1802.04784, ICML 2019.  In the published PDF,
page 4, ``Remarks,'' bullet ``Dependence on $\delta$,'' the authors ask:
\begin{quote}
Whether this robustness to outliers is optimal in the studied setting is an
open question.
\end{quote}
Their Theorems 1--2 and the detailed proof in Appendix A give the upper bound
needed below.

The decisive later source is X. de Juan and S. Mazuelas,
\emph{On the Optimality of the Median-of-Means Estimator under Adversarial
Contamination}, arXiv:2510.07867, NeurIPS 2025.  Its equation (2) and
Theorem 3.2 identify
\[
 \sigma\left(\sqrt{\frac{\log(2/\eta)}{N}}+\sqrt{\alpha}\right)
\]
as the finite-variance minimax order and prove that scalar median-of-means,
with the number of blocks of order $\log(1/\eta)+\alpha N$, attains it.
That paper discusses scalar means, not the MONK min--max estimator.  The
connection below is therefore an implication rather than an explicit answer.

\section*{Rate-level answer}

Let $Z=K(\cdot,X)$ take values in a separable RKHS $\Hcal$, let
$\mu=\mathbb EZ$, and let
$\Sigma=\mathbb E[(Z-\mu)\otimes(Z-\mu)]$ be trace class.  From $N$ clean
i.i.d. observations, an adversary may inspect the sample and replace at most
$m$ entries.  The same convention is used for $m$ contaminated pairs in the
two-sample MMD problem.

\begin{theorem}[MONK is optimal in its outlier dependence]
There are universal constants $c,C>0$ with the following property.  If
$0<\eta\le 1/4$ and
\[
 m+\log(1/\eta)\le cN,
\]
then MONK can be formed with $Q\asymp m+\log(1/\eta)$ blocks so that, with
probability at least $1-\eta$,
\begin{equation}
 \|\widehat\mu_{P,Q}-\mu_P\|_{\Hcal}
 \le C\max\left\{
 \sqrt{\frac{\Tr\Sigma_P}{N}},
 \sqrt{\frac{\|\Sigma_P\|\,[m+\log(1/\eta)]}{N}}
 \right\}.                                      \tag{1}
\end{equation}
Likewise the MONK MMD estimator satisfies
\begin{equation}
 |\widehat{\MMD}_Q(P,R)-\MMD(P,R)|
 \le C\max\left\{
 \sqrt{\frac{\Tr\Sigma_P+\Tr\Sigma_R}{N}},
 \sqrt{\frac{(\|\Sigma_P\|+\|\Sigma_R\|)
 [m+\log(1/\eta)]}{N}}
 \right\}.                                      \tag{2}
\end{equation}
The corruption term in (1), and also that in (2), cannot be improved in
order over the class of distributions with only a finite second moment.  More
precisely, for every $1\le m\le N/4$ and every estimator, there is already a
one-dimensional linear-kernel instance of covariance at most $\sigma^2$ and
an adversary replacing at most $m$ observations for which an error
\[
 c\sigma\sqrt{m/N}                               \tag{3}
\]
occurs with a universal positive probability.
\end{theorem}

Thus the answer to the source question is affirmative in the standard
minimax-rate sense: under the stated trace-class covariance assumption,
MONK has the optimal $\sqrt{m/N}$ outlier dependence.  The clean-sample and
confidence terms were already known to be optimal in the source paper's
discussion.

\section*{Why the implication holds}

\paragraph{Upper bound.}
The detailed proof of Theorem 1 in arXiv:1802.04784 establishes, before its
confidence reparameterization, that whenever
$m\le Q(1/2-\gamma)$,
\begin{equation}
 \|\widehat\mu_{P,Q}-\mu_P\|_{\Hcal}
 \le C_0\max\left\{
 \sqrt{\frac{\|\Sigma_P\|Q}{\gamma N}},
 \frac1\gamma\sqrt{\frac{\Tr\Sigma_P}{N}}
 \right\}                                      \tag{4}
\end{equation}
with probability at least $1-\exp(-Q\gamma^2/72)$.  Choose
$\gamma=1/4$ and an integer
\[
 Q\ge \max\{4m+1,,1152\log(1/\eta)\}
\]
within a universal factor of the right-hand side.  Equation (4) becomes (1).
If $Q$ does not divide $N$, discard fewer than $Q$ observations and use
$N'=Q\lfloor N/Q\rfloor$; under the displayed smallness condition,
$N'\asymp N$, and at most $m$ retained observations are still corrupted.
The Appendix A proof of Theorem 2 gives the identical substitution with the
sums of the two covariance norms and traces, proving (2).

This is exactly the block-count balance isolated for scalar median-of-means
in Theorem 3.2 of arXiv:2510.07867.  The source MONK proof supplies the
dimension-free Hilbert-space lift.

\paragraph{Matching lower bound for mean embeddings.}
It is useful to record a self-contained reduction.  Take
$\mathcal X=\mathbb R$ and the continuous positive-definite linear kernel
$K(x,y)=xy$.  Its RKHS is $\mathbb R$, its feature map is $x$, and its kernel
mean embedding is the ordinary scalar mean.  Put
\[
 \rho=\frac{m}{2N},\qquad
 A=\frac{\sigma}{\sqrt{\rho(1-\rho)}},\qquad
 P_{\pm}=(1-\rho)\delta_0+\rho\delta_{\pm A}.
\]
Both alternatives have variance exactly $\sigma^2$ and means
\[
 \mu_{\pm}=\pm\sigma\sqrt{\frac{\rho}{1-\rho}}.
\]
Let $J$ be the number of nonzero observations.  Since
$\mathbb EJ=N\rho=m/2$, Markov's inequality gives
$\Pr(J\le m)\ge1/2$.  After inspecting the sample, the adversary replaces
all nonzero observations by zero whenever $J\le m$.  On this event the
statistician sees the all-zero data set under either sign.  For any fixed
output on that data set, one of the two signs has error at least
\[
 \sigma\sqrt{\frac{\rho}{1-\rho}}
 \ge \sigma\sqrt{\frac{m}{2N}}.
\]
This proves (3) with failure probability at least $1/2$ for deterministic
estimators; the usual conditioning argument gives a universal constant for
randomized estimators.

\paragraph{Matching lower bound for MMD.}
Fix $R=\delta_0$.  Compare $P_0=\delta_0$ with
$P_1=(1-\rho)\delta_0+\rho\delta_A$, with the same $\rho,A$.  For the linear
kernel,
\[
 \MMD(P_0,R)=0,\qquad
 \MMD(P_1,R)=\sigma\sqrt{\frac{\rho}{1-\rho}}.
\]
The observed data are identically all zero under $P_0$ and, after the same
legal attack, with probability at least $1/2$ under $P_1$.  No estimator can
be within half the displayed separation for both alternatives on that common
observation.  This again yields (3) with universal positive probability.

\section*{Scope and provenance}

The theorem settles the dependence on the sample corruption fraction in the
finite-second-moment model.  It does \emph{not} optimize constants, determine
the exact finite-sample replacement breakdown point, or prove that the
source paper's stated 25\% breakdown value is the largest possible.  Those
are distinct questions from minimax rate optimality.

A bounded search used the exact MONK title and question wording, ``MONK
optimal robustness outliers,'' and ``mean embedding adversarial contamination
minimax finite variance.''  It found arXiv:1812.03523, which proves optimal
contamination exponents under stronger $2+\delta$ moments, and the decisive
arXiv:2510.07867, which proves the exact finite-variance scalar MoM rate.
Because the latter makes the core rate result known and the RKHS conclusion is
a direct identification using the earlier MONK bound, this packet is not a
claim of new mathematics.

\section*{References}

\begin{enumerate}
\item M. Lerasle, Z. Szab\'o, T. Mathieu, and G. Lecu\'e,
\emph{MONK---Outlier-Robust Mean Embedding Estimation by Median-of-Means},
arXiv:1802.04784; Proceedings of ICML 2019, vol. 97, pp. 3782--3793.
\item X. de Juan and S. Mazuelas,
\emph{On the Optimality of the Median-of-Means Estimator under Adversarial
Contamination}, arXiv:2510.07867; NeurIPS 2025.
\item S. Minsker, \emph{Uniform bounds for robust mean estimators},
arXiv:1812.03523 (background comparison only).
\end{enumerate}

\end{document}
